\section{语言建模基础与方法依据}

预训练与任务微调共享自回归框架，通过输入组织和监督范围改变学习任务。本节依次说明字符表示、因果注意力、模块计算与参数优化，为模型实现建立基础。

\subsection{自回归建模与条件生成}

给定词元序列$x_1,\ldots,x_T$，自回归语言模型将联合概率分解为
\begin{equation}
 p_\theta(x_1,\ldots,x_T)=\prod_{t=1}^{T}p_\theta(x_t\mid x_{<t}).
 \label{eq:autoregressive}
\end{equation}
训练阶段使用真实历史词元预测下一词元；推理阶段则使用此前生成的结果继续预测。前者通常称为教师强制训练。两种计算方式的条件序列分别来自真实文本和模型输出，因此训练时的条件预测与推理时的连续生成具有不同的输入分布。

Transformer通过注意力建立序列位置间的联系\cite{transformer}，GPT采用因果自注意力实现自回归生成\cite{gpt2}。将已知上联置于序列前部、下联置于后部，即可用同一模型估计条件概率$p(y\mid x)$；赏析任务也可将任务说明和原文作为前缀，随后生成分析正文。采用仅解码器结构的主要原因是预训练与下游生成能够共享目标和实现框架，并便于在有限计算预算下完成端到端训练。

\subsection{字符表示、嵌入学习与位置编码}

中文文本经规范化后，逐字符映射为整数标识。基础词表共13,232项，包括字符、保留的标点及换行，并设置PAD、UNK、BOS、EOS、UP和DOWN六个特殊词元，分别用于补齐、未知字符、序列起止和上下联区域标记。字符粒度直接对应联长，多字词义则由上下文表示与后续词汇增强模块共同建模。

嵌入层从可训练矩阵$E\in\mathbb{R}^{13232\times256}$中取出相应字符向量，并与位置向量相加：
\begin{equation}
 h_i^{(0)}=E[x_i]+P[i].
 \label{eq:embedding}
\end{equation}
该查表操作等价于独热向量与嵌入矩阵相乘，实际计算直接读取矩阵行。字符嵌入通过自回归交叉熵与其他模型参数联合优化。同一字符进入网络时使用同一嵌入，经过多层上下文计算后，其隐状态可以随语境改变。

本实验采用可学习绝对位置嵌入。位置向量$P[i]$显式编码字符次序，使字符相同而排列不同的序列获得不同表示。长度128的可学习位置表与本实验的短上下文设置相匹配。

\subsection{因果多头注意力与信息可见性}

输入隐状态分别经线性变换得到查询$Q$、键$K$和值$V$。缩放点积注意力为
\begin{equation}
 \operatorname{Attn}(Q,K,V)=\operatorname{softmax}\!\left(\frac{QK^{\mathsf T}}{\sqrt{d_k}}+M\right)V.
 \label{eq:attention}
\end{equation}
查询与键决定读取各位置的权重，值向量承载被聚合的信息。$1/\sqrt{d_k}$控制点积量级；因果掩码$M$将未来位置的分数设为负无穷，使其注意力权重为零。隐藏维256划分为4个注意力头，每头64维，输出拼接后再投影回256维。多头在不同表示子空间中并行计算历史位置的相关性。

因果掩码与损失掩码具有不同作用。前者限制可读取的信息，防止模型利用未来答案；后者选择参与监督的位置。本实验采用右侧补齐：有效查询读取其左侧历史，PAD位置的目标标签从损失中排除。两类掩码分别控制特征计算与目标监督。

\subsection{Pre-LN模块、残差连接与输出层}

每个Transformer模块采用归一化前置结构：先执行LayerNorm与注意力，再经残差相加；之后执行另一组LayerNorm、前馈网络和残差相加。其计算可写为
\begin{align}
 u&=h+\operatorname{MHA}(\operatorname{LN}(h)),\\
 h'&=u+\operatorname{FFN}(\operatorname{LN}(u)).
\end{align}
选择Pre-LN参考其训练稳定性分析\cite{preln}。注意力负责跨位置的信息聚合，前馈网络在各位置独立完成$256\to1024\to256$的非线性映射，激活函数采用GELU。残差连接保留直接的信息与梯度通路，训练时以Dropout进行正则化。

六个模块之后执行最终归一化，并预测下一字符。输入嵌入与输出投影共享权重\cite{tied}，即以$E^{\mathsf T}$将256维隐状态映射至词表空间。这既减少参数，也使字符表示接受输入和输出两端的训练信号。

\subsection{交叉熵监督与参数优化}

设$m_t\in\{0,1\}$表示目标位置是否有效，平均交叉熵为
\begin{equation}
 \mathcal{L}_{\mathrm{CE}}=-\frac{\sum_t m_t\log p_\theta(y_t\mid x,y_{<t})}{\sum_t m_t}.
 \label{eq:ce}
\end{equation}
预训练监督正文和EOS；对联微调仅监督下联及EOS；赏析微调仅监督赏析正文及EOS。前缀参与注意力计算，其表示通过目标段损失接受梯度；外部编码器参数则在词汇增强训练中保持冻结。

优化器采用AdamW\cite{adamw}，结合学习率预热、余弦衰减和梯度裁剪。梯度累计以多个微批次构成一次参数更新，损失按实际有效目标总数归一化，使每个有效目标以相同系数参与优化。BF16用于主要前向计算，参数、优化器状态和交叉熵保持FP32。断点恢复同时加载优化器、计数器和随机状态，维持学习率进程与采样序列的连续性。

\subsection{评价指标定义}\label{sec:metrics}

交叉熵与困惑度反映真实目标在给定真实历史下的可预测性，其中$\mathrm{PPL}=\exp(\mathrm{CE})$。\href{https://github.com/mjpost/sacrebleu}{chrF}衡量生成文本与参考文本的字符$n$元组重合\cite{chrf}，按0至100的量表报告参考文本重合程度。

对联的“形式结构通过率”要求三项同时成立：下联全部由汉字构成且与上联等长；任一相同位置不使用上联原字；任意两位置的重复关系保持一致，即对全部$i<j$满足$(x_i=x_j)\Leftrightarrow(y_i=y_j)$。通过率以指定评测集合的全部上联为分母，空输出与长度错误保留在统计中。全文分别标明1,200条五／七言子集与2,617条全长度集合。

赏析的生成指标均以64条固定输入为分母，且排除程序生成的字音首句。自然结束指模型在词元预算内生成EOS。对正文按步长1提取连续四字符片段，设片段总数为$N$、不同片段数为$U$、最高出现次数为$m$；当$m\geq8$，或$m\geq4$且$1-U/N>0.4$时，记为严重四字符循环。赏析指标统计模型生成正文，程序字音首句单独展示。
